\section{Small-norm descent for the free83 interface}
\label{sec:smallnorm}
To make the nonextension theorem self-contained, we prove the exact
arithmetic lemma imported by its original research packet from
Report43 \cite{smallnorm,freeaudit}. The variable $y_*$ below is the
positive auxiliary coordinate, not the construction exponent $y$.

\begin{lemma}[Small norms and the negative equality orbit]
\label{lem:smallnorm}
Let $H\ge2$, $v\ge0$, $y_*>0$ be integers and
$N_*=Hv^2-(H-1)y_*^2$.
\begin{enumerate}
\item If $N_*<0$, then $N_*\le-(H-1)$.
\item If $0<N_*<H$, then $N_*$ is a positive perfect square.
\item If $N_*=-(H-1)$, every nonnegative-$v$ solution is
\[
 v=2(H-1)\psi_{2H-1}(r),\qquad y_*=\chi_{2H-1}(r),\qquad r\ge0.
\]
\end{enumerate}
\end{lemma}
\begin{proof}
If $N_*<0$, then $y_*>v$. Put $t_*=y_*-v\ge1$; then
\begin{equation}\label{eq:smallidentity}
 -N_*=(H-1)t_*^2+v\bigl(2(H-1)t_*-v\bigr).
\end{equation}
If $v\le2(H-1)t_*$, this is at least $H-1$. Otherwise negativity
implies
\[
 v<\bigl(H-1+\sqrt{H(H-1)}\bigr)t_*<(2H-1)t_*.
\]
The integral transformation
\begin{align*}
 v'&=(2H-1)v-2(H-1)y_*=v-2(H-1)t_*,\\
 y_*'&=(2H-1)y_*-2Hv=(2H-1)t_*-v
\end{align*}
has $0<v'<v$, $y_*'>0$, and preserves $N_*$ by expansion.
Induction on $v$ proves the first assertion.

If $0<N_*<H$, then $v\ge y_*+1$ would give
$N_*=H(v^2-y_*^2)+y_*^2\ge H$, impossible. Thus $y_*\ge v$.
When $y_*=v$ the norm is $v^2$. Otherwise $t_*\ge1$ and
positivity in \eqref{eq:smallidentity} requires
$v>2(H-1)t_*$. Also $v\ge(2H-1)t_*$ would give
$N_*\ge Ht_*^2\ge H$, since the quadratic in $v$ is increasing
beyond its positive root. The same descent has positive coordinates
and smaller $v$, so eventually $y_*=v$; the norm is a square.

For $N_*=-(H-1)$, descend while $v>2(H-1)t_*$.
At a terminal point \eqref{eq:smallidentity} forces $t_*=1$ and
$v=0$ or $2(H-1)$. The latter maps to $(v',y_*')=(0,1)$.
The inverse transformation has matrix
\[
 \begin{pmatrix}2H-1&2(H-1)\\2H&2H-1\end{pmatrix}.
\]
Its iterates on $(0,1)$ are exactly the displayed Pell formula,
as verified by the initial value and recurrence. Conversely every such
iterate has the required norm. This includes $v=0$ at $r=0$.
\end{proof}

\begin{proposition}[Exhaustion of norm signs]\label{prop:signs}
Under \eqref{eq:freenorms}--\eqref{eq:freeproduct} with $A\ge4$,
$f,S,y_*>0$, one has $S\ge2$. If $\Na>0$, then
$\Na=z^2$ for some positive integer $z$ with $z^2\mid\D$.
If $\Na<0$, then necessarily
\begin{equation}\label{eq:negativebranch}
 f=1,\quad S=A,\quad \Ns=-1,\quad \Na=-\D,\quad
 V=\pm2\D\psi_{2A^2-1}(r)\quad(r\ge0).
\end{equation}
\end{proposition}
\begin{proof}
Both norm factors are nonzero divisors of $\D$ and have the same
sign. If $S=1$, then $\Ns=\D f^2-1>1$ is coprime to $\D$,
contradicting divisibility. Set $H=S^2\ge4$.

Suppose $\Na<0$. The small-norm lemma gives
$|\Ns|(H-1)\le\D$, in particular $H-1\le\D$.
If $f\ge2$, the strong equation and $|\Ns|\le\D$ instead give
$H=\D f^2-\Ns\ge3\D$, a contradiction.
Thus $f=1$. Write $b_*=-\Ns\ge1$; then $H=\D+b_*$ and
$b_*(\D+b_*-1)\le\D$, forcing $b_*=1$.
Consequently $S=A$, $\Ns=-1$, $\Na=-\D$.
The equality orbit of Lemma~\ref{lem:smallnorm}, with $v=|V|$,
gives the last formula.

Suppose $\Na>0$; then $\Ns>0$. We show $\Na<H$.
If $f\ge2$, $H\ge3\D>\D/\Ns=\Na$.
If $f=1$, then $\D=H+\Ns$. The value $\Ns=1$ would imply
$A^2-S^2=2$, impossible for positive integers, so $\Ns\ge2$.
Now $\Ns H>H+\Ns=\D$, because $H\ge4$, giving
$H>\D/\Ns=\Na$ again. Lemma~\ref{lem:smallnorm} makes
$\Na$ a positive square, and its divisibility gives $z^2\mid\D$.
\end{proof}

\section{The nonsquarefree Pell lattice and residue separation}
\label{sec:lattice}
A tempting but incorrect shortcut is to assume every strong-$\D$
solution has $f=\chi_A(m)$ and $S=\D\psi_A(m)$. When $\D$
is nonsquarefree this can miss solutions. For example,
$A=26$, $\D=675$ has $f=2,S=45$ satisfying
$675\cdot2^2-45^2=675$; the least positive $\chi_{26}(m)$
is $26$, not $2$. The correct lattice is the following.

\begin{lemma}[Underlying integer Pell unit]\label{lem:lattice}
Let $A\ge4$ be even and $\D=d b_0^2$ with $d$ squarefree.
There exist even $a\ge2$, odd $b\ge1$, and odd $e\ge1$ with
\[
 a^2-db^2=1,\quad
 A+b_0\sqrt d=(a+b\sqrt d)^e.
\]
Put $C_*=\psi_a(e)$ and $\delta_*=a^2-1=db^2$. Then
\begin{equation}\label{eq:latticeidentities}
 b_0=bC_*,\quad \D=\delta_*C_*^2,\quad
 C_*\psi_A(t)=\psi_a(et)\quad(t\ge0).
\end{equation}
Every positive strong-$\D$ solution is exactly
\begin{equation}\label{eq:stronglattice}
 \D f^2-S^2=\D
 \quad\Longleftrightarrow\quad
 f=\chi_a(m),\quad S=\delta_*C_*\psi_a(m),\quad m\ge1.
\end{equation}
Moreover $\gcd(\chi_a(m),b_0)=1$ for every $m\ge0$.
\end{lemma}
\begin{proof}
Since $\D\equiv3\pmod4$, both $b_0$ and $d$ are odd, and
$d\equiv3\pmod4$. Also $d>1$ because $\D$ is not a square.
Among positive integer-coordinate norm-one units for $d$, choose
the least one $u_0=a+b\sqrt d>1$; existence is supplied by
$A+b_0\sqrt d$. Every positive unit is a power of $u_0$:
choose $k$ with $u_0^k\le u<u_0^{k+1}$; the quotient is an
integer-coordinate norm-one unit in $[1,u_0)$, hence is one by
minimality. Its coordinates have the claimed signs because for a
positive real norm-one unit $w\ge1$, the conjugate is $w^{-1}$,
so its two coordinates are nonnegative. This proves generation
without requiring a classification of the full ring of integers.

If $a$ were odd, the norm equation modulo four would force $b$ even,
and all powers would have odd first coordinate, contrary to even $A$.
Thus $a$ is even, $b$ is odd, and modulo two the coefficients of
$u_0^k$ alternate between $(1,0)$ and $(0,1)$. Therefore the exponent
$e$ representing $A+b_0\sqrt d$ is odd. Comparing coordinates and
then powers proves \eqref{eq:latticeidentities}.

If $\D f^2-S^2=\D$, then $b_0^2\mid S^2$, so $b_0\mid S$.
Write $S=b_0s$; then $s^2=d(f^2-1)$, and squarefreeness gives
$s=dv$. Thus $f^2-dv^2=1$, and generation gives
$f=\chi_a(m)$, $v=b\psi_a(m)$ for $m\ge1$ because $S>0$.
This proves \eqref{eq:stronglattice}; its converse follows from the norm.

The norm also implies $\gcd(\chi_a(m),b)=1$. Suppose an odd prime
$\ell_0$ divides both $\chi_a(m)$ and $C_*$.
In the commutative ring
$\mathbb F_{\ell_0}[s]/(s^2-\delta_*)$, the element $u=a+s$
is a unit of norm one. The two divisibilities imply
$u^{2m}=-1$ and $u^{2e}=1$. Raising these identities to $e$ and
$m$, respectively, gives $-1=1$ because $e$ is odd, impossible.
This ring argument remains valid if its quadratic is reducible or has
a repeated root. The prime two does not divide $b_0=bC_*$.
Thus $\gcd(\chi_a(m),b_0)=1$ as claimed.
\end{proof}

\begin{lemma}[Opposite-parity Pell residues]\label{lem:parity}
For $a\ge2$, $m\ge1$, and $F_* =\chi_a(m)$, the sets
\[
 \{\pm\psi_a(t)\bmod F_*:t\ge0\text{ even}\},\qquad
 \{\pm\psi_a(t)\bmod F_*:t\ge0\text{ odd}\}
\]
are disjoint.
\end{lemma}
\begin{proof}
For the formal unit $u=a+\sqrt{a^2-1}$, its coefficientwise
identity modulo $F_*$ is $u^{2m}=-1$; this follows from the
norm and $\chi_a(m)\equiv0$. Reduction modulo $2m$ therefore
changes only a sign. Reflection of a reduced rank $t$ to $2m-t$
preserves the $\psi$ residue, because
$u^{2m-t}=-u^{-t}$ has second coefficient $\psi_a(t)$.
Both operations preserve parity. Hence every residue has a signed
representative $\pm\psi_a(j)$ with $0\le j\le m$ and
$j\equiv t\pmod2$.

For $m=1$ the representatives are zero and one modulo $a\ge2$.
For $m\ge2$ write $b_1=\psi_a(m)$, $b_2=\psi_a(m-1)$.
Strict increase and recurrence give $b_1>(2a-1)b_2$, while
$F_*=ab_1-b_2$. Consequently $b_1+b_2<F_*$:
indeed $F_*-(b_1+b_2)=(a-1)b_1-2b_2>0$.
Distinct-parity representatives have distinct indices, so a nonzero
difference has absolute value less than $F_*$ and a positive sum is
at most $b_1+b_2<F_*$. Neither can be zero modulo $F_*$.
Zero and nonzero representatives are included in this argument.
\end{proof}

\section{Proof of the free83 nonextension theorem}
\label{sec:freeproof}
Assume for contradiction that \eqref{eq:freeV}--\eqref{eq:freeproduct}
hold under Theorem~\ref{thm:free}'s hypotheses.

\subsection{Positive auxiliary norms reduce to one}
By Proposition~\ref{prop:signs}, $\Na=z^2>0$ with
$z^2\mid\D=d b_0^2$. As $d$ is squarefree, $z\mid b_0$.
Multiplying the strong equation by $z^2$ yields
\[
 \D(zf)^2-(zS)^2=\D.
\]
Lemma~\ref{lem:lattice} therefore gives $zf=\chi_a(m)$ for
some $m\ge1$. Its coprimality conclusion forces $z=1$.
Thus every positive branch has
\begin{equation}\label{eq:unitbranch}
 \Na=1,\quad \Ns=\D,\quad
 f=\chi_a(m),\quad S=\delta_*C_*\psi_a(m),\quad m\ge1.
\end{equation}
This is a proof of normalization for this interface, rather than an
assumed generic unit theorem for the full free83 polynomial.

\subsection{The unit branch contradicts rank parity}
By Lemma~\ref{lem:pellclass}, $\Na=1$ implies
\[
 SV=\varepsilon\chi_S(t_*),\quad y_*=\psi_S(t_*),\quad
 t_*\ge1,\quad \varepsilon\in\{+1,-1\}.
\]
Divisibility by $S$ forces $t_*$ odd. Write $t_*=2h+1$.
The strong identity makes $S^2\equiv-\D=1-A^2\pmod f$,
so Lemma~\ref{lem:oddpoly} gives
\[
 V\equiv\varepsilon Q_h(1-A^2)
 =\varepsilon(-1)^h\psi_A(t_*)\pmod f.
\]
The literal expression \eqref{eq:freeV} instead gives
$V\equiv-c=-\psi_A(p)\pmod f$. Multiplying the resulting signed
congruence by $C_*$, and using \eqref{eq:latticeidentities}, yields
\[
 \psi_a(et_*)\equiv\pm\psi_a(ep)\pmod{\chi_a(m)}.
\]
The first rank is odd because $e,t_*$ are odd; the second is even
because $p$ is even. This contradicts Lemma~\ref{lem:parity}.
We multiplied by $C_*$; no unproved division by $C_*$ modulo $f$
is involved. The value of $R$ disappears in this step, so this
positive-branch obstruction does not require $R$ to be odd.

\subsection{The negative branch contradicts the literal parity}
Proposition~\ref{prop:signs} forces the entire negative branch
\eqref{eq:negativebranch}, and in particular makes $V$ even.
But $f=1$ in \eqref{eq:freeV} gives
\[
 V=c(T-1)-R.
\]
The recurrence modulo two gives $\psi_A(p)\equiv p\pmod2$,
so $c$ is even, whereas $R$ is odd. The displayed expression for
$V$ is therefore odd, a contradiction. The equality orbit's
$V=0$ case is also excluded. Positive and negative branches exhaust
the nonzero product $\Na\Ns=\D>0$, proving
Theorem~\ref{thm:free}.

\subsection{Application to the authenticated Report45 tuples}
Those tuples have $A=Y(X+1)+2\equiv2\pmod4$,
$p\equiv0\pmod4$, $c=\psi_A(p)$, $R\equiv3\pmod4$, and
$P_5=1$, exactly as recorded in Section~\ref{sec:setup}.
Hence all satisfy the theorem's hypotheses. The independent free83
audit initially certified the arithmetic theorem conditionally on these
outer invariants because the historical construction directory was
unavailable. The recovered authenticated Report45 construction supplied
here verifies the invariants and closes that source-level application.
It does not extend the conclusion to other outer tuples.

